\section{Experimental Setup}\label{sec:setup}

\subsection{Data and Noise}
All experiments use Fashion-MNIST~\cite{fmnist}, with 60{,}000 training and 10{,}000 test grayscale $28\times28$ images, loaded through \texttt{torchvision}. Pixels are normalized with mean 0.5 and standard deviation 0.5. No validation split is held out, and the test split is used for monitoring and final evaluation. The autoencoder is trained on the full training split. To keep the circuit simulation tractable, the denoisers use only part of the data: the EPQD models use 30 training batches and 8 test batches of latent patches, with at most 8{,}000 patches for quantum fine-tuning, and the MLP and U-Net use 50 training and 10 test batches. Noise is added to latent features, not to pixels, with $\tilde p=p+\epsilon$ and $\epsilon\sim\mathcal{N}(0,0.3^2I)$. Only the standalone stitching network is trained with a smaller noise level of 0.15.

\subsection{Environment and Configuration}
The models are written in PyTorch (2.1 or later), and the circuits in PennyLane (0.40 or later) run on the classical \texttt{default.qubit} simulator, so no quantum hardware or hardware noise is involved. Training uses a GPU when available and otherwise the CPU, and the FID evaluation runs on the CPU. The experiments were run on [hardware description]. All models use Adam with cosine annealing, and Table~\ref{tab:hyper} lists their epochs, learning rates and batch sizes. Each model is trained once with no repetitions, and a fixed seed of 42 is used for the FID evaluation and the V2 benchmark noise. The small differences between the 4-qubit models should therefore not be read as statistically established.

\begin{table}[t]
\centering
\caption{Training settings. Multi-stage models list one value per stage.}
\label{tab:hyper}
\begin{tabular}{lccc}
\toprule
\textbf{Model} & \textbf{Epochs} & \textbf{Learning rate} & \textbf{Batch}\\
\midrule
Conv-AE & 20 & $10^{-3}$ & 128\\
StitchingCNN & 15 & $10^{-3}$ & 128\\
MLP, U-Net & 20 & $10^{-3}$ & 128\\
EPQD (warm-up, fine-tune) & 10, 8 & $10^{-3}$, $3\times10^{-3}$ & 128\\
EPQD-6Q-V2 (3 phases) & 15, 8, 6 & $10^{-3}$, $2\times10^{-3}$, $5\times10^{-4}$ & 128, 128, 16\\
\bottomrule
\end{tabular}
\end{table}

\subsection{Baselines and Metrics}
The four 4-qubit EPQD models and EPQD-6Q-V2 are compared with the patch-wise MLP and the U-Net. All of them receive the same noisy latent maps from the same frozen autoencoder and use MSE as the training loss. The decoded clean latent map is the best possible result, and the decoded noisy latent map is the lower reference.

Images are decoded and mapped to $[0,1]$ before evaluation. We report the mean squared error, $\mathrm{MSE}=\frac{1}{N}\sum_i(x_i-\hat x_i)^2$, and the peak signal-to-noise ratio, $\mathrm{PSNR}=10\log_{10}(1/\mathrm{MSE})$. We also report the structural similarity index (SSIM), which compares local luminance, contrast and structure. The Fr\'echet Inception Distance (FID) compares the Inception-V3 feature distributions of real and reconstructed images,
\begin{equation}
    \operatorname{FID}=\|\mu_r-\mu_g\|_2^2+
    \operatorname{Tr}\!\left(\Sigma_r+\Sigma_g-2(\Sigma_r\Sigma_g)^{1/2}\right),
\end{equation}
and is computed on 1{,}000 test images. The V2 benchmark uses 64 test images and also reports RMSE, mean absolute error and cosine similarity. Lower MSE and FID and higher PSNR and SSIM indicate better results.
